% Mise en forme du sommaire du dossier de projet.
%
% Avant-propos, chapitres et annexes sont tous au même rang dans le document :
% le sommaire tient donc sur un seul niveau, aligné sur la marge, plutôt que
% d'indenter les chapitres sous l'avant-propos comme s'ils en dépendaient.
%
% On redéfinit directement les commandes internes de LaTeX plutôt que de passer
% par tocloft, absent du cache local et qui entrerait de toute façon en conflit
% avec titlesec, déjà chargé par le gabarit.

\makeatletter

% Entrées de sommaire : pas de retrait, pas de réserve pour un numéro
% (la numérotation « 1. », « 2. »… fait partie du titre), et un peu d'air
% entre les lignes pour aérer une liste de vingt entrées.
\renewcommand*\l@subsection[2]{%
  \vskip 3\p@
  \@dottedtocline{2}{0pt}{0pt}{#1}{#2}%
}

% Points de conduite plus serrés que les 4.5 mu par défaut.
\renewcommand{\@dotsep}{2}

% Colonne des numéros de page, et marge droite réservée pour l'accueillir.
\renewcommand{\@pnumwidth}{2em}
\renewcommand{\@tocrmarg}{2.8em}

\makeatother

% Les noms de classes, de fichiers et les routes composés en chasse fixe ne se
% coupent pas : sans marge de manœuvre, TeX les laisse déborder dans la marge.
% On lui autorise un étirement supplémentaire des espaces sur les paragraphes
% qu'il ne parvient pas à justifier autrement.
\setlength{\emergencystretch}{3.5em}


% --- Pagination -------------------------------------------------------------

% Pandoc compose les images en figures flottantes sans indication de placement,
% ce qui interdit à LaTeX de les poser à l'endroit où elles sont écrites : il ne
% lui reste que le haut ou le bas de page, et les figures en attente s'accumulent
% jusqu'à remplir des pages entières, voire à en laisser une blanche. On rétablit
% le placement « ici » et on autorise une figure à occuper une plus grande part
% de la page, pour qu'elle reste près du paragraphe qui l'annonce.
\floatplacement{figure}{htbp}
\renewcommand{\topfraction}{0.92}
\renewcommand{\bottomfraction}{0.75}
\renewcommand{\textfraction}{0.06}
\renewcommand{\floatpagefraction}{0.7}
\setcounter{topnumber}{3}
\setcounter{bottomnumber}{2}
\setcounter{totalnumber}{4}

% Une ligne seule en haut ou en bas de page se lit mal : mieux vaut laisser la
% page un peu plus courte et emmener le paragraphe entier sur la suivante.
\widowpenalty=10000
\clubpenalty=10000
\displaywidowpenalty=10000

% Corollaire du point précédent : on laisse les bas de page respirer plutôt que
% d'étirer les blancs entre paragraphes pour aligner artificiellement la dernière
% ligne sur la marge basse.
\raggedbottom

% Les planches de l'annexe E sont annoncées par une ligne de titre suivie d'une
% ou deux images. \besoindeplace, placé avant le titre, exige la hauteur donnée :
% s'il ne reste pas assez de page, le titre passe sur la suivante avec ses images
% au lieu de rester seul en bas.
\usepackage{needspace}
\newcommand{\besoindeplace}[1]{\Needspace*{#1}}
